% Title:  Computing the adjoint matrix of multivariate polynomial matrices given 
%         by straight-line programs

 
% Authors:  B. Casta�o, J. Heintz, J. Llovet


% Affiliation:  B. Casta�o, J. Llovet:
%               Department of Mathematics, University of Alcal�, Madrid, Spain
%
%               J. Heintz:
%               Department of Mathematics, Faculty of Science, University of Cantabria,
%               Santander, Spain
%               and
%               Department of Mathematics, University of Buenos Aires, Buenos Aires,
%               Argentina


% Abstract (in LaTeX): 

\documentstyle[12pt]{article}
\newcounter{cms}
\pagestyle{myheadings}       % estilo de p\'agina
\setlength{\parskip}{0.0cm}  % espacio entre parrafos 
\setlength{\parsep}{0.0cm}  
%\renewcommand{\itemsep}
% ---------------- Margenes --------------------------
\oddsidemargin  0.7 cm  % margen + 1 pulgada  0.7 cm
\evensidemargin 0.7 cm  % margen + 1 pulgada  0.7 cm
\textwidth       15.5 cm  % anchura del texto   15  cm
\textheight      23 cm  % longitud del texto  21  cm
%----------------- Cabecera --------------------------
\topmargin   0 pt  % distancia a la cabacera
\headheight  .5 cm  % anchura de la cabacera
\headsep     0 cm  % desde la cabecera hasta el texto
\topskip     0 pt  % borde del texto a primera linea
\setlength{\partopsep}{-10 pt} % separacion entre texto y primer item
%-------------- la � ---------------------------------
\newcommand{\gn}{\~{n}}
% ------------- Distancia entre l�neas -----------------
\renewcommand{\baselinestretch}{1.0}
\pagestyle{empty}

\begin{document} % <<<<<<<<<<<<<<<<<<<<<<<<<<<<<<<<<<<<<<<<<<<<<<

\begin{center}
{\Large {\bf Abstract}}
\end{center}
\vspace{.3 cm}
\setlength{\parskip}{0 cm}

In this paper we describe a new implementation of an algorithm for computing the adjoint matrix of a multivariate 
polynomial matrix. The innovation of our approach consists on using an arithmetic circuit implemented in the computer as 
a directed acyclic graph (DAG) for representing the polynomials, the matrices and the algorithms as well. The mentioned 
DAG represents a straight-line program (SLP) when executing an evaluation algorithm, otherwise known as pebble game 
across the arithmetic circuit. 

The spoken of procedure has been described extensively in previously published papers, presented at ACA [{\it Llovet 
(1998), Casta\~no (2000)}]. In all these papers we have described very interesting results. Thus, it was shown that by 
using above mentioned SLP representation it was possible to compute the characteristic polynomial of multivariate 
polynomial matrices and the greatest common divisor of two multivariate polynomials in less time. It also required less 
computer memory than other, more common computer algebra systems such as MAPLE or MAGMA.

In the current paper we describe a design that is based on the same mathematical principle. Subsequently, we have 
managed to compute the adjoint matrix of a $n\times n$ square polynomial matrix $A=(A_{ij}= P_{ij}(x_1,\ldots,x_m))$ by 
deriving the determinant of $A$ given in DAG representation with respect to all its entries. 
This entirely novel approach, which required backward-derivation we made possible by relying on Baur-Strassen's 
algorithm in the Jacques Morgenstern's version[{\it Morgenstern}]. 

{\bf References. }
\renewcommand{\labelenumi}{[\arabic{enumi}]}

\begin{enumerate}
\setlength{\itemsep}{0 pt}
\setlength{\parsep}{2 pt}
\item J. Llovet, B. Casta\gn o, R. Mart\'{\i}nez, Computing the characteristic polynomial of multivariate polynomial 
matrices given by straight-line programs, Mathematics and Computers in Simulation 45, (1998), pp. 39-57.
\item B. Casta\gn o, J. Heintz, J. Llovet, R. Mart\'{\i}nez, On the data structure straight-line program and its 
implementation in symbolic computation, Mathematics and Computers in Simulation 51, (2000), pp. 497-528.
\item J. Morgenstern, How to compute fast a function and all its derivatives A variation on the theorem of 
Baur-Strassen, preprint.
\end{enumerate}
\end{document}
